% This is part of Mes notes de mathématique
% Copyright (c) 2011-2019
%   Laurent Claessens
% See the file fdl-1.3.txt for copying conditions.

\label{CHAPGroupes}

Pour rappel, la notion de groupe est définie en \ref{DEFooBMUZooLAfbeM}.

%+++++++++++++++++++++++++++++++++++++++++++++++++++++++++++++++++++++++++++++++++++++++++++++++++++++++++++++++++++++++++++ 
\section{Somme à valeurs dans un groupe abélien}
%+++++++++++++++++++++++++++++++++++++++++++++++++++++++++++++++++++++++++++++++++++++++++++++++++++++++++++++++++++++++++++

Si \( S\) est un ensemble fini, nous savons de la proposition \ref{PROPooJLGKooDCcnWi} qu'il existe un unique \( N\in \eN\) pour lequel il existe une bijection \( \varphi\colon \{ 0,\ldots, N \}\to S\). Cette bijection n'est a priori pas unique.

\begin{lemmaDef}[\cite{MonCerveau}]       \label{DEFooLNEXooYMQjRo}
    Soient un groupe abélien \( (G,+)\) ainsi qu'un ensemble fini \( I\) contenant \( n\) éléments. Soit une application \( f\colon I\to G \). Si \( \sigma_1,\sigma_2\colon \{1,\ldots, n \}\to I\) sont deux bijections, alors\footnote{Pour rappel, le symbole \( \sum_{i=1}^n\) est défini par \ref{DEFooNEVNooJlmJOC}.}
    \begin{equation}
        \sum_{i=1}^nf\big( \sigma_1(i) \big)=\sum_{i=1}^nf\big( \sigma_2(i) \big).
    \end{equation}
    La valeur commune est notée
    \begin{equation}
        \sum_{i\in I}f(i)
    \end{equation}
\end{lemmaDef}

La définition \ref{DEFooLNEXooYMQjRo} donne lieu à un certain nombre de remarques.
\begin{enumerate}
    \item
        Elle donne la somme sur un ensemble fini. Un problème avec les ensembles infinis (outre la convergence) est l'ordre de sommation. Si vous voulez sommer sur \( \eZ\), dans quel ordre le faire ?
    \item
        Pour aller plus loin, et sommer sur des ensembles infinis, il faut regarder la définition \ref{DefHYgkkA}. 
\end{enumerate}

\begin{proposition}     \label{PROPooJBQVooNqWErk}
    Soient un groupe abélien \( (G,+)\), un ensemble fini \( I\) , une application \( f\colon I\to G\) et une bijection \( \sigma\colon I\to I\). Alors
    \begin{equation}
        \sum_{i\in I}f(i)=\sum_{i\in I}f\big( \sigma(i) \big).
    \end{equation}
\end{proposition}

Si nous avons une application \( L\colon S\to S\), nous notons
\begin{equation}
    \sum_{s\in S}f\big( L(s) \big)=\sum_{s\in S}(f\circ L)(s).
\end{equation}
Cette façon d'écrire donne une interprétation pour la notation \( \sum_{g\in G}f(hg)\) qui arrive dans la proposition \ref{PROPooWJQQooFINSEc}. Il s'agit de considérer l'application \( L_h\) du lemme \ref{LEMooBIBFooBHxFYC}, de considérer\footnote{Le fait que \( L_h\) soit une bijection n'a pas d'importance ici.}
\begin{equation}        \label{EQooQQBEooFDOBVG}
    \sum_{g\in G}f(hg)=\sum_{g\in G}(f\circ L_h)(g)
\end{equation}
et de faire tourner la définition \ref{DEFooLNEXooYMQjRo}. La même chose tient pour définir \( \sum_{g\in G}(gh)\) à l'aide de \( R_h\).

\begin{proposition}[\cite{MonCerveau}]     \label{PROPooWJQQooFINSEc}
    Soient un groupe fini \( G\) et une fonction \( f\colon G\to A\) où \( A\) est un anneau commutatif. Alors
    \begin{equation}
        \sum_{g\in G}f(g)=\sum_{g\in G}f(gh)=\sum_{g\in G}f(hg)
    \end{equation}
    pour tout \( h\in G\).
\end{proposition}

\begin{proof}
    Nous avons une bijection \( \varphi\colon \{ 0,\ldots,  N \}\to G\) garantie par la proposition \ref{PROPooJLGKooDCcnWi}. La définition est que
    \begin{equation}
        \sum_{g\in G}f(g)=\sum_{i=0}^Nf\big( \varphi(i) \big).
    \end{equation}
    Par ailleurs, le lemme \ref{LEMooBIBFooBHxFYC} donne une bijection \( L_h\colon G\to G\) et permet de considérer la composée
    \begin{equation}
        \begin{aligned}
            \varphi'\colon \{ 0,\ldots,  N \}&\to G \\
            \varphi'=L_h\circ \varphi.
        \end{aligned}
    \end{equation}
    La proposition \ref{DEFooLNEXooYMQjRo} nous permet d'utiliser la bijection \( \varphi'\) au lieu de \( \varphi\) pour exprimer la somme \( \sum_{g\in G}\). Ensuite un jeu de notation utilisant \eqref{EQooQQBEooFDOBVG} donne
    \begin{equation}
        \begin{aligned}[]
            &\sum_{g\in G}f(g)=\sum_{i=0}^Nf\big( \varphi(i) \big)=\sum_{i=0}^Nf\big( \varphi'(i) \big)=\sum_{i=0}^N(f\circ L_h\circ \varphi)(i)\\
            &\quad=\sum_{i=0}^N(f\circ L_h)\big( \varphi(i) \big)=\sum_{g\in G}(f\circ L_h)(g)=\sum_{g\in G}f(hg).
        \end{aligned}
    \end{equation}
    En ce qui concerne \( \sum_{g\in G}f(gh)\), c'est la même chose, en utilisant \( R_h\) au lieu de \( L_h\).
\end{proof}

Tout cela nous permet de faire une somme sympathique et bien connue.
\begin{lemma}
    Soit \( n\in \eN\). Nous avons
    \begin{equation}
        \sum_{k=0}^nk=\frac{ n(n+1) }{ 2 }.
    \end{equation}
\end{lemma}

\begin{proof}
    La preuve est pratiquement immédiate par récurrence. Nous allons donner une preuve plus «constructive», qui formalise l'idée classique d'écrire la somme à l'endroit et à l'envers.


    Nous notons \( S\) la somme \( \sum_{k=0}^nk\). Le lemme \ref{DEFooLNEXooYMQjRo} dit que si \( \sigma_i\colon \{ 0,\ldots, n \}\to \{ 0,\ldots, n \}\) sont deux bijections, alors \( \sum_{k=0}^nf\big( \sigma_1(k) \big)=\sum_{k=0}^nf\big( \sigma_2(k) \big)\). Nous sommes intéressé au cas \( f(i)=i\).

    En prenant \( \sigma_1(k)=k\) et \( \sigma_2(k)=n-k\), nous avons
    \begin{equation}
        S=\sum_{k=0}^nk=\sum_{k=0}^n(n-k).
    \end{equation}
    Donc
    \begin{equation}
        2S=\sum_{k=0}^n\big( k+(n-k) \big)=\sum_{k=0}^nn=n\sum_{k=0}^n1=n(n+1).
    \end{equation}
    En divisant par deux, nous obtenons le résultat annoncé.
\end{proof}

Dans le même ordre d'idée, pour le produit.
\begin{proposition}     \label{PROPooQMUDooQQVRIe}
    Si \( E\) est un ensemble fini et si \( G\) est un groupe abélien, alors pour toute fonction \( f\colon E\to G\) et pour toute permutation \( \sigma\) de \( E\),
    \begin{equation}
        \prod_{i\in E}f(i)=\prod_{i\in E}f\big( \sigma(i) \big)
    \end{equation}
\end{proposition}

%+++++++++++++++++++++++++++++++++++++++++++++++++++++++++++++++++++++++++++++++++++++++++++++++++++++++++++++++++++++++++++
\section{Du vocabulaire dans les groupes}
%+++++++++++++++++++++++++++++++++++++++++++++++++++++++++++++++++++++++++++++++++++++++++++++++++++++++++++++++++++++++++++

%---------------------------------------------------------
\subsection{Centre, centralisateur, normalisateur}
%---------------------------------------------------------

\begin{definition}[\cite{Kropholler}]\label{defGroupeCentre}
Soit \( G\) un groupe, et \( H\) un sous-ensemble de \( G \).
\begin{enumerate}
\item
Le \defe{centre}{centre!d'un groupe} d'un groupe \( G\) est l'ensemble des éléments de \( G\) qui commutent avec tous les autres:
\begin{equation}
    Z_G=\{ g\in G\tq gh=hg\;\forall g\in G \}.
\end{equation}

\item
Si \( h\in G\) nous notons \( Z_G(h)\) le \defe{centralisateur}{centralisateur} de \( h\) dans \( G\) :
\begin{equation}
    Z_G(h)=\{g\in G\tq gh = hg\}.
\end{equation}
C'est l'ensemble des éléments de \( G\) qui commutent avec \( g\).

\item    Le \defe{centralisateur}{centralisateur} de \( H\subset G\) est
    \begin{equation}
        Z_G(H)=\{ g\in G\tq gh=hg\,\forall h\in H\};
    \end{equation}
    il contient donc tous les éléments de \( G\) qui commutent avec ceux de \( H\).
\end{enumerate}
\end{definition}

\begin{definition}\label{DEFGroupeNormalisateur}
        Si \( H\) est un sous-groupe, son \defe{normalisateur}{normalisateur} est
    \begin{equation}
        N_G(H)=\{ g\in G\tq gH=Hg \}.
    \end{equation}
\end{definition}

\begin{remark}
Les définitions précedentes ne sont utiles que si le groupe \(G \) n'est pas abélien~\footnote{Définition~\ref{DEFGroupeAbelien}.}. Si \( G \) est abélien, alors clairement~\footnote{Vérifiez-le!}, \(Z_G = G,\ Z_G(g) = G\) pour tout \(g \in G,\ Z_G(H) = G,\ H_G(H) = G\), et tous les sous-groupes sont normaux.
\end{remark}

\begin{remark}
Les définitions précédentes ont un certain lien entre elles; en effet:
  \begin{enumerate}
    \item dans la définition~\ref{DEFGroupeCentre}, le centre est le centralisateur de \(G \), et le centralisateur d'un ensemble \( H \) est l'intersection des centralisateurs des éléments \( h \) de \( H \);
    \item un sous-groupe \( H \) est normal si et seulement si \( N_G(H) = G \).
  \end{enumerate}
\end{remark}


%+++++++++++++++++++++++++++++++++++++++++++++++++++++++++++++++++++++++++++++++++++++++++++++++++++++++++++++++++++++++++++
\section{Action de groupes}
%+++++++++++++++++++++++++++++++++++++++++++++++++++++++++++++++++++++++++++++++++++++++++++++++++++++++++++++++++++++++++++

\begin{definition}[Thème~\ref{THEMEooKZHBooRCULcr}]  \label{DefActionGroupe}
    Une \defe{action de groupe}{action}\index{action} \( G\) sur un ensemble \( E\) est la donnée, pour chaque élément \( g \in G\), d'une fonction \(\phi_g : E \to E \), de telle sorte que:
    \begin{gather*}
        \phi_{e}(x) = x, \hspace{2em} \forall x \in E;\\
        \phi_{gh}(x) = \phi_g (\phi_h (x)),  \hspace{2em} \forall g,h \in G, \forall x \in E.
     \end{gather*}
     On dit dans ce cas que \( G \) \defe{agit}{action} sur \( E \).
\end{definition}

\begin{lemma}
    Pour tout \( g\in G\),
    \begin{enumerate}
        \item
            L'application \( \phi_g\colon E\to E\) est injective,
        \item
            Pour l'inverse : \( (\phi_g)^{-1}=\phi_{g^{-1}}\).
    \end{enumerate}
\end{lemma}

\begin{proof}
    Si \( x,y\in E\) sont tels que \( \phi_g(x)=\phi_g(y)\) alors en appliquant \( \phi_{g^{-1}}\) aux deux membres nous trouvons
    \begin{equation}
        (\phi_{g^{-1}}\phi_g)(x)=(\phi_{g^{-1}}\phi_g)(y),
    \end{equation}
    ce qui donne \( x=y\) parce que \( \phi_{g^{-1}}\phi_g=\phi_{g^{-1}g}=\phi_e=\id\).

    Les dernières trois égalités écrites disent que \( \phi_{g^{-1}}\) est l'inverse\footnote{Si vous décidez de dire ça a un jury dans un concours, soyez prêts à préciser les domaines.} de \( \phi_g\).
\end{proof}

Pour alléger les notations, on convient d'écrire $g \cdot x$, voire plus simplement $gx$ au lieu de \( \phi_g(x) \). Le deuxième axiome d'action de groupe dit que la notation $ghx$ ne souffre d'aucune ambiguïté.

\begin{definition}[Quelques notions autour de l'action]

    Si \( G\) agit sur un ensemble \( E\), nous notons \( G\cdot x\) l'\defe{orbite}{orbite!d'un point sous une action} de \( x\in E\) sous l'action de $G$:
\begin{equation*}
    G\cdot x = \{ gx \tq g \in G\}.
\end{equation*}

Nous notons \( G_x\) ou \( \Fix(x)\) le \defe{stabilisateur}{stabilisateur} de \( x\) :
\begin{equation}
    \Fix(x)=G_x=\{ g\in G\tq g\cdot x=x \}.
\end{equation}

    Pour \( g\in G\), nous notons enfin \( \Fix(g)\) le \defe{fixateur}{fixateur} de \( g\) :
\begin{equation}
    \Fix(g)=\{ x\in E\tq g\cdot x=x \}.
\end{equation}

\end{definition}

\begin{definition}  \label{DefuyYJRh}
    L'action de \( G\) sur \( E\) est \defe{fidèle}{fidèle (action)}\index{action!fidèle} si l'identité est le seul élément de \( G\) à fixer tous les points de \( E\), c'est-à-dire si \( gx=x\,\forall x\in E\Rightarrow g=e\).
\end{definition}

Un exemple d'action fidèle tout à fait non trivial sera donné avec l'action du groupe modulaire sur le plan de Poincaré dans le théorème~\ref{ThoItqXCm}.

Le groupe \( G\) agit toujours sur lui même à gauche et à droite. L'action à gauche est \( g\cdot h=gh\); celle à droite est \( g\cdot h=hg^{-1}\).

\begin{definition}      \label{DEFooCORTooEeOLPT}
    L'action \defe{adjointe}{action!adjointe} définie par \( g\cdot h=ghg^{-1}\) est une manière pour un groupe d'agir sur lui-même par automorphismes. Cela est souvent noté \( \AD(g)h=ghg^{-1}\).
\end{definition}
En effet pour tout \( g\in G\), l'application \( \AD(g)\colon G\to G\) est un automorphisme de \( G\).

Si \( H\) est un sous-groupe de  \( G\), nous notons \( G/H\) le quotient de $G$ par la relation \( g\sim gh\) pour tout \( h\in H\). Lorsque la distinction est importante, nous noterons \( (G/H)_g\)\nomenclature[R]{$(G/H)_g$}{classes à gauche} pour les classes à gauche et \( (G/H)_d\) pour les classes à droite.

Nous avons une relation d'équivalence à gauche et une à droite. D'abord
\begin{equation}
    x\sim_g y\Leftrightarrow xh=y
\end{equation}
pour un certain \( h\in H\). Ensuite
\begin{equation}
    x\sim_d y\Leftrightarrow hx=y
\end{equation}
pour un certain \( h\in H\).

Le lemme suivant est une généralisation du théorème de Lagrange~\ref{ThoLagrange}.

\begin{lemma}
    L'ensemble \( (G/H)_g\) est fini si et seulement si l'ensemble \( (G/H)_d\) est fini. S'il en est ainsi, alors \( (G/H)_g\) et \( (G/H)_d\) ont même cardinal qui vaut l'indice de \( H\) dans \( G\).
\end{lemma}

\begin{proof}
    L'application
    \begin{equation}
        \begin{aligned}
            f\colon (G/H)_g&\to (G/H)_d \\
            [x]_g&\mapsto [x^{-1}]_d
        \end{aligned}
    \end{equation}
    est une bijection bien définie. En effet si \( x\sim_g y\), nous avons \( h\in H\) tel que \( y^{-1}h=x^{-1}\), c'est-à-dire que \( x^{-1}\sim_d y^{-1}\) et \( f\) est bien définie. Le fait que \( f\) soit surjective est évident. Pour l'injectivité, soient \( x, y \in G \) tels que
    \begin{equation}
        f([x]_g)=f([y]_g).
    \end{equation}
    Alors \( x^{-1}\sim_d y^{-1}\), ce qui implique l'existence de \( h\in H\) tel que \( hx^{-1}=y^{-1}\), ou encore que \( xh^{-1}=y\), ce qui signifie que \( x\sim_gy\).

    Pour l'énoncé à propos de l'indice, nous procédons en plusieurs étapes simples.
    \begin{enumerate}
        \item
            Les classes (les éléments de \( (G/H)_g\)) forment une partition de $G$.
        \item
            Toutes les classes ont le même nombre d'éléments par la bijection
            \begin{equation}
                \begin{aligned}
                    f\colon [x]_g&\to [y]_g \\
                    xh&\mapsto yh.
                \end{aligned}
            \end{equation}
        \item
            Le nombre d'éléments dans une classe est égal à \( | H |\) par la bijection
            \begin{equation}
                \begin{aligned}
                    g\colon [x]_g&\to H \\
                    xh&\mapsto h.
                \end{aligned}
            \end{equation}
    \end{enumerate}
    Par conséquent
    \begin{equation}
        | G |=| H |\cdot \text{nombre de classes}=| H |\cdot\text{cardinal de $(G/H)_g$},
    \end{equation}
    et nous avons bien
    \begin{equation}
        \text{cardinal de }(G/H)_g=\frac{ | G | }{ | H | }=| G:H |.
    \end{equation}
\end{proof}

\begin{proposition}[Orbite-stabilisateur\cite{Combes}]     \label{Propszymlr}
    Soit \( G\) un groupe agissant sur un ensemble \( E\) et \( x\in E\).
    \begin{enumerate}
        \item
            Les ensembles \( G\cdot x\) et \( G/G_x\) sont équipotents.
        \item       \label{ITEMooCWUGooCOFHYk}
            L'orbite de \(\Fix(x)\) est finie si et seulement si \( \Fix(x)\) est d'indice fini dans \( G\). Dans ce cas nous avons
            \begin{equation}        \label{EqnLCHCE}
                \Card(G\cdot x)=| G:\Fix(x) |.
            \end{equation}
            Une autre façon d'écrire la même formule :
            \begin{equation}        \label{EqCewSXT}
                | G |=| \Fix(x) | |\mO_x |.
            \end{equation}
    \end{enumerate}
\end{proposition}
\index{équation!orbite-stabilisateur}
C'est la formule \eqref{EqnLCHCE} qui est nommée \wikipedia{fr}{Action_de_groupe_(mathématiques)\#Formule_des_classes.2C_formule_de_Burnside}{formule des classes} sur wikipédia.

\begin{proof}
    \begin{enumerate}
        \item
    Soit l'application
    \begin{equation}
        \begin{aligned}
            \psi\colon G\cdot x&\to G/G_x \\
            a\cdot x&\mapsto [a].
        \end{aligned}
    \end{equation}
    Cette application est bien définie parce que si \( a\cdot x=b\cdot x\), alors il existe \( h\in G_x\) tel que \( b=ah\), et par conséquent \( [a]=[b]\). Cette application est une bijection et par conséquent \( G\cdot x\) est équipotent à \( G/G_x\).
    \item
        Soit \( y\in \mO_x\) et \( A_y=\{ g\in G\tq g\cdot x=y \}\). L'ensemble \( A_y\) est une classe à gauche de \( \Fix(x)\), par conséquent \( | A_y |=|\Fix(x)|\) pour tout \( y\in\mO_x\). Les \( A_y\) pour différents \( y\) sont disjoints et nous avons de plus
        \begin{equation}
            \bigcup_{y\in\mO_x}A_y=G.
        \end{equation}
        Les ensembles \( A_y\) divisent donc \( G\) en \( | \mO_x |\) paquets de \( | \Fix(x) |\) éléments. D'où la formule \eqref{EqCewSXT}.

    \end{enumerate}
\end{proof}

\begin{corollary}       \label{CORooRRVHooTyCjZZ}
    Soit \( C_g\) la classe de conjugaison d'un élément  \( g\) du groupe fini \( G\). Alors
    \begin{equation}
        \Card(C_g)=| G:Z_G(g) |
    \end{equation}
    où $Z_G(g)$ est le centralisateur de \( g\) dans \( G\)\footnote{Définition~\ref{defGroupeCentre}.} de \( G\).
\end{corollary}

\begin{proof}
    Cela est une application de la proposition~\ref{Propszymlr} (formule \eqref{EqnLCHCE}) dans le cas de l'action adjointe de \( G\) sur lui-même.

    En effet, si nous considérons l'action adjointe, l'orbite est la classe de conjugaison : \( C_g=G\cdot g\). Et le stabilisateur de \( g\) pour l'action adjointe n'est autre que le centralisateur de \( g\) :
    \begin{subequations}
        \begin{align}
        \Fix(g)&=\{ h\in G\tq h\cdot g=g \}\\
        &= \{ h\in G\tq hgh^{-1}=g \}\\
        &=\{ h\in G\tq gh=hg \}\\
        &=Z_G(g).
        \end{align}
    \end{subequations}

    Donc la formule \( \Card(G\cdot g)=| G:G_g |\) devient, dans le cas de l'action adjointe de \( G\) sur lui-même : \( \Card(C_g)=| G:Z_G(g) |\).
\end{proof}

\begin{lemma}
    Soit \( G\) un groupe agissant sur l'ensemble \( E\). On définit \( x\sim x'\) si et seulement s'il existe \( g\in G\) tel que \( g\cdot x=x'\). Alors
    \begin{enumerate}
        \item
            la relation \( \sim\) est une relation d'équivalence.
        \item
            la classe \( [x]\) est l'orbite \( \mO_x\) de \( x\) sous \( G\).
    \end{enumerate}
\end{lemma}

\begin{corollary}[Équation des orbites]\index{équation!des orbites} \label{CorARFVMP}
    Soit \( G\) un groupe agissant sur l'ensemble \( E\) et \( \mO_1,\ldots, \mO_k  \) la liste des orbites (distinctes). Alors
    \begin{enumerate}
        \item
            \( E=\bigcup_i\mO_i\), l'union est disjointe,
        \item
            \( \Card(E)=\sum_i\Card(\mO_i)\).
    \end{enumerate}
\end{corollary}

\begin{definition}  \label{DefcSuYxz}
    Soit \( G\) un groupe agissant sur l'ensemble \( E\). Un \defe{domaine fondamental}{domaine!fondamental d'une action}\index{fondamental!domaine d'une action}\index{action!domaine fondamental} ou une \defe{transversale}{transversale} est une partie de \( E\) contenant un et un seul élément de chaque orbite.
\end{definition}
Autrement dit, les images des éléments d'un domaine fondamental forment une partition de l'ensemble :
\begin{equation}
    E=\bigsqcup_{g\in G}g(F)
\end{equation}
où  \( g(F) = \phi_g(F) = \{ \phi_g(x) \tq x \in F\} \). L'union est disjointe, c'est-à-dire que si \( g\neq g'\), alors \( g(F)\cap g'(F)=\emptyset\).


%TODO:
% - définir ce qu'est un p-groupe, au bon endroit (à côté du théorème de Lagrange?).
% - pour l'instant, la définition de p-groupe est là où on l'utilise, c'est-à-dire un peu avant les théorèmes de Sylow.
% - Il vaut peut-être mieux créer un thème ``indice dans un groupe'' dans lequel on parlera entre autres du théorème de Lagrange et de p-groupe.


\begin{proposition}[Équation des classes\cite{FabricegPSFinis}]     \label{PropUyLPdp}
    Soit \( G\), un groupe fini opérant sur un ensemble \( E\). Si \( E''\) est un ensemble contenant exactement un élément de chaque orbite dans \( E\setminus\Fix_G(E)\), alors
    \begin{equation}        \label{EqobuzfK}
        | G |=| \Fix_G(E) |+\sum_{x\in E''}\frac{ | G | }{ | \Fix_G(x) | }.
    \end{equation}
    Si de plus \( G\) est un $p$-groupe, alors
    \begin{equation}    \label{EqbzLEVJ}
        | E |=| \Fix_G(E) |\mod p.
    \end{equation}
\end{proposition}

\begin{proof}
    Par le corolaire~\ref{CorARFVMP}, nous avons \( | G |=\sum_{x\in E'}| \mO_x |\) où \( E'\) est une transversale.  En séparant la somme entre les orbites à un élément et les autres,
    \begin{equation}    \label{EqeggkBs}
        | G |=\Card(\Fix_G(E))+\sum_{x\in E''}\frac{ | G | }{ | \Fix_G(x) | }
    \end{equation}  \label{EqDgYbhm}
    où nous avons utilisé le fait que \( | G |=| \Fix_G(x) | |\mO_x |\).

    Si \( G\) est un \( p\)-groupe alors si \( x\in E''\), \( \Fix_G(x)\) est un sous-groupe propre de \( G\) et donc \( | \Fix_G(x) |\) est un diviseur propre de \( | G |\). Du coup la fraction \( | G |/|\Fix_G(x)|\) est une puissance non nulle de \( p\) et l'équation \eqref{EqobuzfK} devient immédiatement \eqref{EqbzLEVJ}.
\end{proof}


\begin{corollary}[Équation des classes]\index{équation!des classes}
    Soit \( G\), un groupe et \( C_1\),\ldots, \( C_l\) la liste de ses classes de conjugaison contenant plus de un élément. Alors
    \begin{equation}        \label{EqkgGmoq}
        \Card(G)=\Card\big( Z(G) \big)+\sum_i| G:Z_{g_i} |=\Card\big( Z(G) \big)+\sum_i\frac{ \Card(G) }{ \Card\big( \Fix(g_i) \big) }
    \end{equation}
    si \( g_i\in C_i\).
\end{corollary}

\begin{proof}
    Étant donné que les classes de conjugaison sont disjointes, le cardinal de \( G\) est bien la somme des cardinaux de ses classes. Les classes ne contenant que un seul élément sont celles des éléments de \( Z(G)\). En ce qui concerne les autres orbites, \( \Card(C_{g_i})=| G:Z_{g_i} |\) par le théorème orbite-stabilisateur (proposition~\ref{Propszymlr}).
\end{proof}



\begin{theorem}[Théorème de Burnside\cite{FabricegPSFinis}] \label{ThoImkljy}
    Le centre d'un \( p\)-groupe non trivial est non trivial.
\end{theorem}

\begin{proof}
    Soit \( G\) un $p$-groupe non trivial. Nous considérons l'action adjointe \( G\) sur lui-même. Les points fixes de cette action sont les éléments du centre :
    \begin{equation}
        \mZ_G=\{ z\in G\tq \sigma_x(z)=z\forall x\in G \}=\Fix_G(G).
    \end{equation}
    Nous utilisons l'équation aux classes \eqref{PropUyLPdp} pour dire que \( | G |=| \mZ_G |\mod p\). Mais \( | \mZ_G |\) n'est pas vide parce qu'il contient l'identité. Donc \( | \mZ_G |\) est au moins d'ordre \( p\).
\end{proof}

\begin{proposition} \label{PropssttFK}
    Si \( p\) est un nombre premier, tout groupe d'ordre \( p\) ou \( p^2\) est abélien.
\end{proposition}
Rappel : un groupe d'ordre \( p\) ou \( p^2\) est automatiquement un $p$-groupe.

\begin{proof}
    Si \( | G |=p\), alors le théorème de Cauchy~\ref{ThoCauchyGpFini} nous donne l'existence d'un élément d'ordre \( p\). Cet élément est alors automatiquement générateur, \( G\) est cyclique et donc abélien.

    Si par contre \( G\) est d'ordre \( p^2\), alors les choses se compliquent (un peu). D'après le théorème de Burnside~\ref{ThoImkljy}, le centre \( \mZ\) n'est pas trivial; il est alors d'ordre \( p\) ou \( p^2\). Supposons qu'il soit d'ordre \( p\) et prenons \( x\in G\setminus\mZ\). Alors le stabilisateur de \( x\) pour l'action adjointe contient au moins \( \mZ\) et \( x\), c'est-à-dire que \( |\Fix_G(x)|\geq p+1\). Étant donné que \( \Fix_G(x)\) est un sous-groupe, son ordre est automatiquement \( 1\), \( p\) ou \( p^2\). En l'occurrence, il doit être \( p^2\) (parce que plus grand que \( p\)), et donc \( x\) doit être central, ce qui est une contradiction.
\end{proof}


\begin{theorem}[\wikipedia{fr}{Action_de_groupe_(mathématiques)}{Formule de Burnside}]      \label{THOooEFDMooDfosOw}
    Si \( G\) est un groupe fini agissant sur l'ensemble fini \( E\) et si \( \Omega\) est l'ensemble des orbites, alors
    \begin{equation}    \label{EqTUsblv}
        \Card(\Omega)=\frac{1}{ | G | }\sum_{g\in G}\Card\big( \Fix(g) \big).
    \end{equation}
\end{theorem}
\index{Burnisde!formule}
\index{formule!Burnside}

\begin{proof}
    Nous considérons l'ensemble
    \begin{equation}
        A=\{ (g,x)\in G\times E\tq gx=x \},
    \end{equation}
    et nous en calculons le cardinal de deux façons. D'abord
    \begin{subequations}
        \begin{align}
            \Card(A)&=\sum_{x\in E}\Card\{ g\in g\tq gx=x \}\\
            &=\sum_{x\in E}\Card(\Fix(x))\\
            &=\sum_{\omega\in \Omega}\sum_{x\in \omega}\Card(\Fix(x))\\
            &=\sum_{\omega\in \Omega}\frac{ | G | }{ \Card(\omega) }     \label{EqyVtkyf}\\
            &=| G |.
        \end{align}
    \end{subequations}
    Pour obtenir \eqref{EqyVtkyf} nous avons utilisé l'équation des classes \eqref{EqCewSXT}. L'autre façon de calculer \( \Card(A)\) est de regrouper ainsi :
    \begin{equation}
        \Card(A)=\sum_{g\in G}\Card\{ x\in E\tq gx=x \}=\sum_{g\in G}\Card(\Fix(g)).
    \end{equation}
    En égalisant les deux expressions de \( \Card(A)\) nous trouvons
    \begin{equation}
        | G |\Card(\Omega)=\sum_{g\in G}\Card(\Fix(g)).
    \end{equation}
\end{proof}

\begin{proposition}
    Soit \( G\) un groupe et \( H\), un sous-groupe du centre de \( G\).
    \begin{enumerate}
        \item
            \( H\) est normal dans \( G\).
        \item
            Si \( G/H\) est monogène, alors \( G\) est abélien.
        \item
            Si \( G\) est fini de centre \( Z\), alors \( | G:H |\) n'est pas premier.
    \end{enumerate}
\end{proposition}

% TODO: preuve

\begin{theorem}     \label{THOooRGSTooIWyhqt}
    Soit \( G\) un groupe cyclique\footnote{Définition \ref{DefHFJWooFxkzCF}.} d'ordre \( n\).
    \begin{enumerate}
        \item
            Tout sous-groupe de \( G\) est cyclique.
        \item
            Pour chaque diviseur \( d\) de \( n\), il existe un unique sous-groupe \( H_d\) de \( G\) d'ordre \( d\).
    \end{enumerate}
    Si \( a\) est un générateur de \( G\), alors \( H_d\) peut être décrit des façons suivantes :
    \begin{equation}
        H_d=\{ x\in G\tq x^d=e \}=\{ x\in G\tq\exists y\in G\tq y^{n/d}=x \}=\langle a^{n/d}\rangle.
    \end{equation}
\end{theorem}

% TODO: preuve

\begin{definition}      \label{DEFooQDHPooCfDEuL}
    Soit \( G\) un groupe agissant sur un ensemble \( E\). Nous disons que l'action est \defe{transitive}{transitive}\index{action!transitive} si elle possède une seule orbite. L'action est \defe{libre}{libre!action}\index{action!libre} si \( g\cdot x=g'\cdot x\) implique \( g=g'\).
\end{definition}


